% Part: first-order-logic
% Chapter: model-theory
% Section: overspill

\documentclass[../../../include/open-logic-section]{subfiles}

\begin{document}

\olfileid[tr]{mod}{bas}{ove}
\olsection{Taşma}

\begin{thm}
\ollabel{overspill} Bir tümceler kümesi~$\Gamma$ gelişigüzel büyük
sonlu modellere sahipse sonsuz bir modele de sahiptir.
\end{thm}

\begin{proof}
$\Gamma$ dilini sayılabilir sonsuz sayıda yeni $c_0,c_1,\dots$
sabitini ekleyerek genişletin ve
$\Gamma\cup\{c_i\neq c_j:i\neq j\}$ kümesini ele alın. $\Gamma$'nın
gelişigüzel büyük sonlu modellere sahip olması, her $m>0$ için
$n\geq m$ olacak biçimde $\Gamma$'nın $n$ kardinaliteli bir modelinin
bulunması demektir. Bu, $\Gamma\cup\{c_i\neq c_j:i\neq j\}$ kümesinin
sonlu sağlanabilir olduğunu gerektirir. Tıkızlık gereği bu kümenin bir
$\Struct M$ modeli vardır. Bütün $c_i\neq c_j$ eşitsizliklerini
sağladığı için bu modelin evreni sonsuz olmak zorundadır.
\end{proof}

\begin{prop}
\ollabel{inf-not-fo}
Herhangi bir birinci dereceden dilde, !!a{structure}~$\Struct M$
içinde ancak ve ancak !!{structure} ifadesinin evreni~$\Domain M$
sonsuz olduğunda doğru olan hiçbir $!A$ tümcesi yoktur.
\end{prop}

\begin{proof}
Böyle bir $!A$ bulunsaydı $\lnot !A$ tam ve yalnızca sonlu
!!{structure}s içinde doğru olurdu. Dolayısıyla gelişigüzel büyük sonlu
modellere sahip olduğu halde sonsuz bir modeli bulunmazdı; bu da
\olref{overspill} ile çelişir.
\end{proof}

\end{document}
